A linear probe reads a clinical finding from the visual representation that a vision-language model's language model consumes, and activation steering suggests the converse: writing the probe direction back should move the model's answer about the finding, making the direction a handle on the concept. Current evaluations score encoding and question answering, leaving the handle untested. We test it with \emph{ownership}: at one steering strength on the same patients, does a concept's own direction move its own answer more than the directions fitted for every competing finding? Across \cfCheckpoints{} checkpoints on NIH ChestX-ray14, CheXpert Plus and COCO, clinical concepts are readable in \cfChestReadablePct{}\% of chest checkpoint--concept cells and answerable in \cfChestAnswerablePct{}\%, but owned in \cfChestOwnedPct{}\%: most clinical writes do not move their own answer beyond random directions of the same steering strength, and the same directions move the other findings' answers. Under the same procedure, \cfCocoOwnedPct{}\% of natural-image object cells are owned, and the gap persists among cells the model answers well (\cfStratChestTopPct{}\% against \cfStratNatTopPct{}\%). Non-clinical attributes of the radiographs share the deficit, and checkpoints that share a vision tower, and therefore write identical vectors, own different concepts: the deficit follows the chest-radiograph representation and the model that reads it. A direction fitted to the model's own answers provides the handle that the label direction lacks, and on chest radiographs the two directions are nearly orthogonal.
